\section{El momento AlexNet: Software 2.0 pasa de un avance de investigación a una señal industrial}

\subsection{Por qué 2012 se convirtió en un punto de inflexión}

La entrevista retoma el resultado contundente de AlexNet en la competencia ImageNet de 2012. Más importante aún, Jensen describió el contexto interno de NVIDIA en ese momento: la empresa también trabajaba en visión por computadora, pero los métodos tradicionales tenían el progreso estancado, y el avance de AlexNet equivalió a «dar justo en el blanco» de ese punto de dolor interno.

\begin{figure}[H]
\centering
\includegraphics[width=0.88\textwidth]{frame-alexnet-slide.jpg}
\caption{La entrevista muestra la portada del artículo de AlexNet, subrayando su significado paradigmático \protect\footnotemark}
\end{figure}
\footnotetext{Intervalo de tiempo en el video: 00:11:20--00:11:30.}

\begin{knowledgebox}{La migración de «programar» a «entrenar»}
La definición de Jensen sobre esta migración puede resumirse así:
\begin{itemize}
    \item Software 1.0: las personas escriben las reglas y la máquina las ejecuta.
    \item Software 2.0: las personas proporcionan los datos y la máquina aprende las reglas.
\end{itemize}

Cuando la escala y la complejidad del problema superan el límite de lo que se puede mantener con reglas escritas a mano, la ruta 2.0 erosiona de manera constante el territorio de la 1.0. Esta es también la razón por la que la posterior explosión de los LLM pudo propagarse rápidamente entre industrias.
\end{knowledgebox}

\subsection{«Reasoned hope» y la segunda confirmación}

Jensen describió el juicio de aquel año como «reasoned hope»: ni un optimismo sin evidencia, ni una espera conservadora de validación, sino una apuesta de alto riesgo y alta recompensa dentro de las restricciones de la ingeniería y del mercado. La aparición de AlexNet le dio a NVIDIA una segunda confirmación:
\begin{itemize}
    \item La ruta del hardware paralelo iba en la dirección correcta.
    \item Las cargas de trabajo de entrenamiento se convertirían en la corriente principal a largo plazo.
    \item El ecosistema de software debía diseñarse alrededor de sistemas de aprendizaje y no de algoritmos fijos.
\end{itemize}

\begin{warningbox}{Error común: AlexNet fue solo «un modelo que ganó una competencia»}
Ver a AlexNet como una victoria de benchmark hace perder el mensaje central. Lo que en realidad se validó fue la regularidad conjunta de «datos + cómputo + modelo escalable», así como la posición de infraestructura de la GPU dentro de esa regularidad.
\end{warningbox}

\subsection{Resumen de la sección}
Para NVIDIA, el significado de AlexNet no fue «subirse a una tendencia», sino la confirmación de que la ruta de CUDA en la que había apostado una década antes podía sostener el siguiente paradigma de cómputo.

